\section{Лекция 9. Дифференциальные уравнения 1-го порядка, не разрешённые относительно производной}
\noeqref{eq:l9_imp_ode, eq:l9_ex_parab, eq:l9_cauchy_imp, eq:l9_ift_eq}
\noeqref{eq:l9_ift_deriv, eq:l9_exist_cond, eq:l9_res_y, eq:l9_res_y_diff}
\noeqref{eq:l9_res_y_ans, eq:l9_res_x, eq:l9_res_x_diff, eq:l9_res_x_ans}
\noeqref{eq:l9_f267, eq:l9_f272}

\begin{flushright}
    26.10.26
\end{flushright}

\subsection{Постановка задачи Коши и геометрический смысл единственности}

В предыдущих лекциях мы изучали дифференциальные уравнения 1-го порядка,
разрешённые относительно производной: $y' = f(x, y)$. Перейдём теперь к
общему случаю уравнений 1-го порядка, \textdef{не разрешённых относительно
производной}:
\begin{equation} \label{eq:l9_imp_ode}
    f(x, y, y') = 0, \tag{1}
\end{equation}
где $f(x, y, p)$ — непрерывная функция трёх переменных в некоторой области
пространства $\R^3$.

\begin{example} ~
    Примеры уравнений, не разрешённых относительно производной:
    \begin{align*}
        (y')^2 - 2 x y' + 4 y &= 0, \\
        (y')^7 - x y &= 0, \\
        (y')^3 - \sin y' &= y.
    \end{align*}
\end{example}

Для уравнения $y' = f(x, y)$, разрешённого относительно производной, задача
Коши задавалась одним начальным условием $y(x_0) = y_0$, поскольку в каждой
точке $(x_0, y_0)$ наклон поля направлений $y'(x_0) = f(x_0, y_0)$ определялся
однозначно. Для уравнения \eqref{eq:l9_imp_ode} ситуация принципиально меняется:
алгебраическое уравнение $f(x_0, y_0, y') = 0$ может иметь несколько различных
корней относительно $y'$, то есть в одной точке $(x_0, y_0)$ возникает сразу
несколько допустимых направлений касательных.

\begin{example}
    Рассмотрим уравнение, распадающееся на два линейных множителя:
    \begin{equation} \label{eq:l9_ex_parab}
        (y')^2 - x^2 = 0 \iff (y' - x)(y' + x) = 0. \tag{2}
    \end{equation}

    \begin{proof}[Анализ семейств решений]~\\
        \begin{quote}
            Уравнение \eqref{eq:l9_ex_parab} распадается на совокупность двух
            разрешённых относительно производной уравнений:
            \[
                \left[
                \begin{aligned}
                    y' &= x, \\
                    y' &= -x
                \end{aligned}
                \right.
                \implies
                \left[
                \begin{aligned}
                    y &= \frac{x^2}{2} + C_1, \\
                    y &= -\frac{x^2}{2} + C_2,
                \end{aligned}
                \right.
                \quad C_1, C_2 \in \R.
            \]
            Через каждую точку $(x_0, y_0) \in \R^2$ проходят сразу две
            параболы: одна ветвями вверх ($y = \frac{x^2}{2} + C_1$), а другая
            ветвями вниз ($y = -\frac{x^2}{2} + C_2$). Исследуем поведение
            этих решений в зависимости от выбора начальной точки $(x_0, y_0)$:
            \begin{itemize}
                \item \textbf{Точки вне оси ординат ($x_0 \neq 0$):} В такой
                точке пересекающиеся параболы имеют \textbf{разные производные}:
                $y'(x_0) = x_0$ и $y'(x_0) = -x_0 \neq x_0$. Поэтому, если
                дополнительно зафиксировать начальный наклон касательной
                $y'(x_0) = y_1$ (где $y_1 = x_0$ или $y_1 = -x_0$), то из двух
                парабол однозначно выделится ровно одна интегральная кривая.
                \item \textbf{Точки на оси ординат ($x_0 = 0$):} В любой точке
                $(0, y_0)$ параболы $y = \frac{x^2}{2} + y_0$ и
                $y = -\frac{x^2}{2} + y_0$ касаются друг друга и имеют
                \textbf{одинаковые производные} $y'(0) = 0$. В этих точках
                даже при фиксированной производной $y'(0) = 0$ единственность
                нарушается.
            \end{itemize}
        \end{quote}
    \end{proof}
\end{example}

\begin{figure}[H]
    \centering
    \begin{tikzpicture}[scale=1.15]
        % Оси координат
        \draw[->] (-3.2, 0) -- (3.2, 0) node[right] {$x$};
        \draw[->] (0, -2.2) -- (0, 3.2) node[above] {$y$};

        % Семейство парабол ветвями вверх (оранжевые/красные)
        \foreach \c in {0, 0.9, 1.8} {
            \draw[red!85!orange, thick, domain=-2.1:2.1, samples=60]
                plot (\x, {0.5*\x*\x + \c});
        }

        % Семейство парабол ветвями вниз (зелёные/бирюзовые)
        \foreach \c in {0, 0.9, 1.8} {
            \draw[teal, thick, domain=-2.1:2.1, samples=60]
                plot (\x, {-0.5*\x*\x + \c});
        }

        % Точки неединственности на оси Oy (производные равны: y'(0) = 0)
        \foreach \y in {0, 0.9, 1.8} {
            \fill[red] (0, \y) circle (2pt);
        }
        \node[red, align=left, font=\footnotesize] at (-2.1, 2.55)
            {точки неединственности\\[-1pt] (производные равны)};
        \draw[->, red, thin] (-1.2, 2.3) .. controls (-0.5, 2.0) .. (-0.1, 1.82);
        \draw[->, red, thin] (-1.2, 2.3) .. controls (-0.6, 1.3) .. (-0.1, 0.92);

        % Точки единственности вне оси Oy (производные разные)
        \fill[teal!80!black] (0.9487, 1.35) circle (2pt);
        \fill[teal!80!black] (1.3416, 0.9) circle (2pt);
        \node[teal!80!black, align=left, font=\footnotesize] at (2.4, 2.1)
            {точки единственности\\[-1pt] (производные разные)};
        \draw[->, teal!80!black, thin] (1.55, 1.95) -- (1.02, 1.42);
        \draw[->, teal!80!black, thin] (1.75, 1.75) -- (1.40, 0.98);
    \end{tikzpicture}
    \caption{Интегральные кривые уравнения $(y')^2 - x^2 = 0$: при $x_0 \neq 0$
    кривые двух семейств пересекаются под ненулевым углом (точки
    единственности), а на прямой $x_0 = 0$ касаются друг друга (точки
    неединственности).}
\end{figure}

Приведённый пример мотивирует дать корректную постановку задачи Коши для уравнения
\eqref{eq:l9_imp_ode}.

\begin{definition}
    \textdef{Задачей Коши} для дифференциального уравнения 1-го порядка, не
    разрешённого относительно производной, называется задача нахождения
    дифференцируемой функции $y = y(x)$, удовлетворяющей уравнению и двум
    начальным условиям (значению функции и значению её первой производной в
    начальной точке):
    \begin{equation} \label{eq:l9_cauchy_imp}
        \begin{cases}
            f(x, y, y') = 0, \\
            y(x_0) = y_0, \\
            y'(x_0) = y_1,
        \end{cases} \tag{3}
    \end{equation}
    где тройка чисел $(x_0, y_0, y_1) \in \R^3$ обращает уравнение
    в верное числовое равенство:
    $f(x_0, y_0, y_1) = 0$.
\end{definition}

\subsection{Теорема существования и единственности}

Чтобы свести уравнение \eqref{eq:l9_imp_ode} к уже изученному разрешённому виду
$y' = \phi(x, y)$, нам потребуется фундаментальная теорема математического
анализа о неявной функции.

\begin{example}
    Рассмотрим нелинейное алгебраическое соотношение на плоскости:
    \[
        x y + \cos\left(\frac{x}{y}\right) = 0.
    \]
    Выразить из него переменную $y = f(x)$ в виде явной элементарной формулы
    невозможно, однако при выполнении определённых условий гладкости такое
    функциональное соответствие $y = f(x)$ существует локально в окрестности
    точки $(x_0, y_0)$, удовлетворяющей уравнению.
\end{example}

\begin{theorem}[\textbf{Теорема о неявной функции}]
    Пусть функция $F(x, y)$ определена и непрерывно дифференцируема в
    некоторой окрестности $W \subset \R^2$ точки $(x_0, y_0)$, то есть
    $F \in C^1(W)$, и выполняются два условия:
    \begin{equation} \label{eq:l9_ift_eq}
        F(x_0, y_0) = 0, \qquad
        F'_y(x_0, y_0) = \frac{\partial F}{\partial y}(x_0, y_0) \neq 0. \tag{4}
    \end{equation}
    Тогда существует прямоугольная окрестность
    $U_\eps(x_0) \times V_\delta(y_0) \subset W$ и существует единственная
    непрерывно дифференцируемая функция $y = f(x) \in C^1(U_\eps(x_0))$,
    принимающая значения в $V_\delta(y_0)$, такая что:
    \[
        f(x_0) = y_0 \quad \text{и} \quad
        \forall x \in U_\eps(x_0) : F\bigl(x, f(x)\bigr) \equiv 0.
    \]
    При этом производная неявной функции вычисляется по формуле:
    \begin{equation} \label{eq:l9_ift_deriv}
        f'(x) = -\frac{F'_x\bigl(x, f(x)\bigr)}{F'_y\bigl(x, f(x)\bigr)}. \tag{5}
    \end{equation}
\end{theorem}

\begin{proof}[Доказательство]~\\
    \begin{quote}
        Предположим для определённости, что $F'_y(x_0, y_0) > 0$. Так как
        частные производные $F'_x, F'_y \in C(W)$, найдётся замкнутый
        прямоугольник $\Pi = [x_0 - \alpha, x_0 + \alpha] \times
        [y_0 - \delta, y_0 + \delta] \subset W$, во всех точках которого
        выполняется строгое неравенство $F'_y(x, y) > 0$.

        \textbf{1. Существование и единственность значения $y = f(x)$.}\\
        Поскольку при фиксированном $x = x_0$ функция одной переменной
        $g(y) = F(x_0, y)$ строго возрастает на $[y_0 - \delta, y_0 + \delta]$
        и $F(x_0, y_0) = 0$, на концах отрезка она имеет разные знаки:
        \[
            F(x_0, y_0 - \delta) < 0 < F(x_0, y_0 + \delta).
        \]
        По непрерывности функции $F(x, y)$ найдётся такое $\eps \in (0, \alpha]$,
        что для всех $x \in U_\eps(x_0) = (x_0 - \eps, x_0 + \eps)$ знаки на
        верхней и нижней крышках сохраняются:
        \[
            F(x, y_0 - \delta) < 0 \quad \text{и} \quad F(x, y_0 + \delta) > 0.
        \]
        Тогда при каждом фиксированном $x \in U_\eps(x_0)$ непрерывная функция
        $y \mapsto F(x, y)$ по теореме Больцано — Коши о промежуточном значении
        принимает нулевое значение в некоторой точке $y \in (y_0 - \delta, y_0 + \delta)$,
        а в силу строгого возрастания ($F'_y > 0$) такой корень $y = f(x)$
        \textbf{единственен}. В частности, при $x = x_0$ в силу единственности
        получаем $f(x_0) = y_0$.

        \textbf{2. Непрерывная дифференцируемость функции $f(x)$.}\\
        Дадим аргументу $x \in U_\eps(x_0)$ малое приращение $\Delta x$, которому
        соответствует приращение $\Delta y = f(x + \Delta x) - f(x)$. Так как
        обе точки лежат на кривой $F = 0$, имеем $F(x + \Delta x, y + \Delta y)
        - F(x, y) = 0$. По формуле конечных приращений Лагранжа:
        \[
            F'_x(x + \theta_1 \Delta x, y + \theta_2 \Delta y) \Delta x +
            F'_y(x + \theta_1 \Delta x, y + \theta_2 \Delta y) \Delta y = 0,
            \quad \theta_1, \theta_2 \in (0, 1).
        \]
        Отсюда при $\Delta x \to 0$ в силу непрерывности $F'_x, F'_y$ и
        условия $F'_y \neq 0$ получаем $\Delta y \to 0$ (непрерывность $f(x)$)
        и предельное соотношение \eqref{eq:l9_ift_deriv}:
        \[
            f'(x) = \lim_{\Delta x \to 0} \frac{\Delta y}{\Delta x} =
            -\frac{F'_x(x, f(x))}{F'_y(x, f(x))} \in C(U_\eps(x_0)).
        \]
    \end{quote}
\end{proof}

\begin{figure}[H]
    \centering
    \begin{tikzpicture}[scale=1.15]
        % Оси координат
        \draw[->] (-0.5, 0) -- (4.2, 0) node[right] {$x$};
        \draw[->] (0, -0.5) -- (0, 3.2) node[above] {$y$};

        % Замкнутая линия уровня F(x, y) = 0 (эллипс)
        \draw[blue, thick, dashed] (2.0, 1.5) ellipse (1.5 and 1.0);
        \node[blue, above right] at (3.1, 2.2) {$F(x, y) = 0$};

        % Окрестность точки (x_0, y_0), где F'_y \neq 0
        \def\xO{2.0}
        \def\yO{2.5}
        \draw[red, thick, dashed] (\xO - 0.8, \yO - 0.35)
            rectangle (\xO + 0.8, \yO + 0.35);
        \draw[red, very thick, domain=1.2:2.8, samples=60]
            plot (\x, {1.5 + sqrt(1.0 - ((\x - 2.0)/1.5)^2)});
        \fill[red] (\xO, \yO) circle (2pt)
            node[above] {$(x_0, y_0): F'_y \neq 0$};

        % Проекция окрестности U_\eps(x_0) на ось Ox
        \draw[dashed, gray] (\xO - 0.8, 0) -- (\xO - 0.8, \yO - 0.35);
        \draw[dashed, gray] (\xO + 0.8, 0) -- (\xO + 0.8, \yO - 0.35);
        \draw[<->, red, thick] (\xO - 0.8, -0.25) -- (\xO + 0.8, -0.25)
            node[midway, below] {\footnotesize $U_\eps(x_0)$};

        % Особая точка с вертикальной касательной, где F'_y = 0
        \fill[teal!80!black] (3.5, 1.5) circle (2pt);
        \draw[teal!80!black, thick] (3.5, 0.9) -- (3.5, 2.1);
        \node[teal!80!black, right, font=\footnotesize, align=left]
            at (3.55, 1.5) {$F'_y = 0$\\ (нет однозначной $y(x)$)};
    \end{tikzpicture}
    \caption{Геометрический смысл теоремы о неявной функции: в окрестности
    точки $(x_0, y_0)$, где $F'_y(x_0, y_0) \neq 0$ (касательная не
    вертикальна), кривая $F(x, y) = 0$ однозначно проектируется на ось $Ox$,
    задавая гладкую функцию $y = f(x)$.}
\end{figure}

Применим теперь теорему о неявной функции к исследованию задачи Коши
\eqref{eq:l9_cauchy_imp}.

\begin{theorem}[\textbf{О существовании и единственности решения задачи Коши}]
    Пусть функция $f(x, y, p)$ определена и непрерывно дифференцируема в
    некоторой окрестности $W \subset \R^3$ точки $(x_0, y_0, y_1)$, то есть
    $f \in C^1(W)$, и в самой начальной точке выполняются условия:
    \begin{equation} \label{eq:l9_exist_cond}
        \begin{cases}
            f(x_0, y_0, y_1) = 0, \\
            f'_{y'}(x_0, y_0, y_1) =
            \dfrac{\partial f}{\partial y'}(x_0, y_0, y_1) \neq 0.
        \end{cases} \tag{6}
    \end{equation}
    Тогда в некоторой окрестности $U_\eps(x_0) = (x_0 - \eps, x_0 + \eps)$
    существует \textdef{единственное} решение $y = y(x) \in C^2(U_\eps(x_0))$
    задачи Коши \eqref{eq:l9_cauchy_imp}.
\end{theorem}

\begin{proof}[Доказательство]~\\
    \begin{quote}
        Обозначим $y' = p$ и рассмотрим уравнение $f(x, y, p) = 0$ в
        окрестности точки $(x_0, y_0, y_1) \in \R^3$.

        \textbf{Шаг 1. Локальное разрешение относительно производной $p = y'$.}\\
        По условию $f \in C^1(W)$, $f(x_0, y_0, y_1) = 0$ и
        $\frac{\partial f}{\partial p}(x_0, y_0, y_1) \neq 0$. Согласно теореме
        о неявной функции (Теорема 9.1, применённая к функции трёх переменных),
        существует окрестность $V \subset \R^2$ точки $(x_0, y_0)$ и
        интервал $I_\delta = (y_1 - \delta, y_1 + \delta)$, такие что в
        цилиндре $V \times I_\delta$ уравнение $f(x, y, p) = 0$ равносильно
        явному соотношению
        \[
            p = \phi(x, y),
        \]
        где $\phi \in C^1(V)$ — единственная непрерывно дифференцируемая
        функция, удовлетворяющая начальному условию $\phi(x_0, y_0) = y_1$.

        \textbf{Шаг 2. Применение теоремы Пикара.}\\
        Таким образом, для кривых, у которых $(x, y) \in V$ и
        $y' \in I_\delta$, исходная задача Коши \eqref{eq:l9_cauchy_imp}
        эквивалентна классической задаче Коши для уравнения, разрешённого
        относительно производной:
        \[
            \begin{cases}
                y' = \phi(x, y), \\
                y(x_0) = y_0,
            \end{cases}
        \]
        причём второе начальное условие $y'(x_0) = \phi(x_0, y_0) = y_1$
        выполняется автоматически. Так как $\phi \in C^1(V)$, её частная
        производная
        \[
            \frac{\partial \phi}{\partial y}(x, y) =
            -\frac{f'_y\bigl(x, y, \phi(x, y)\bigr)}{f'_{y'}\bigl(x, y, \phi(x, y)\bigr)}
        \]
        непрерывна в $V$. Следовательно, в некотором замкнутом прямоугольнике
        $\Pi \subset V$ с центром в точке $(x_0, y_0)$ функция $\phi(x, y)$
        непрерывна и удовлетворяет условию Липшица по переменной $y$ (см.
        Замечание 3.1 из Лекции 3). По теореме Пикара (Теорема 3.1) существует
        окрестность $U_\eps(x_0)$, на которой эта задача Коши имеет
        единственное решение $y = y(x)$.
    \end{quote}
\end{proof}

\begin{remark}
    Вернёмся к примеру \eqref{eq:l9_ex_parab}: $f(x, y, y') = (y')^2 - x^2 = 0$.
    Для него частная производная по $y'$ равна $f'_{y'} = 2 y'$. Условие
    $f'_{y'} = 0$ даёт $y' = 0$, откуда из самого уравнения следует $x = 0$.
    Именно поэтому на оси ординат $x = 0$ условие \eqref{eq:l9_exist_cond}
    теоремы 9.2 нарушается, и там возникают точки неединственности!
\end{remark}

\subsection{Методы интегрирования уравнений, не разрешённых относительно производной (метод введения параметра)}

Рассмотрим аналитический метод интегрирования уравнений $f(x, y, y') = 0$,
которые можно явно разрешить относительно одной из переменных — $y$ или $x$.
Основная идея состоит во введении параметра $p = y'$, при котором переменные
$x, y$ и параметр $p$ связаны дифференциальным соотношением:
\[
    dy = p \, dx.
\]

\subsubsection{Случай 1. Уравнение разрешено относительно $y$: $y = F(x, y')$}

Пусть дифференциальное уравнение приведено к виду $y = F(x, y')$, где
$F \in C^1$. Положив $y' = p$, запишем систему:
\begin{equation} \label{eq:l9_res_y}
    \begin{cases}
        y = F(x, p), \\
        dy = p \, dx.
    \end{cases} \tag{7}
\end{equation}

Возьмём полный дифференциал от обеих частей первого уравнения системы
\eqref{eq:l9_res_y}:
\[
    dy = \frac{\partial F}{\partial x} \, dx +
    \frac{\partial F}{\partial p} \, dp.
\]
Заменяя в левой части дифференциал $dy$ на $p \, dx$, получаем дифференциальное
уравнение 1-го порядка, связывающее переменные $x$ и $p$:
\begin{equation} \label{eq:l9_res_y_diff}
    p \, dx = \frac{\partial F}{\partial x} \, dx +
    \frac{\partial F}{\partial p} \, dp \iff
    \left(p - \frac{\partial F}{\partial x}(x, p)\right) dx =
    \frac{\partial F}{\partial p}(x, p) \, dp. \tag{8}
\end{equation}

\begin{enumerate}
    \item Интегрируя уравнение \eqref{eq:l9_res_y_diff}, находим зависимость
    $x = x(p, C)$ и подставляем её во второе соотношение $y = F(x, p)$. Общее
    решение записывается в параметрической форме:
    \begin{equation} \label{eq:l9_res_y_ans}
        \begin{cases}
            x = x(p, C), \\
            y = F\bigl(x(p, C), p\bigr).
        \end{cases} \tag{9}
    \end{equation}
    \item Если при решении уравнения \eqref{eq:l9_res_y_diff} выполнялось
    деление на множитель, зависящий от $p$ (например, при $dp = 0 \implies
    p = p_0 = \text{const}$), то корни $p = p_0$ уравнения
    $p_0 - \frac{\partial F}{\partial x}(x, p_0) = 0$ проверяются отдельно: они
    могут давать особые решения $y = F(x, p_0)$.
\end{enumerate}

\subsubsection{Случай 2. Уравнение разрешено относительно $x$: $x = F(y, y')$}

Пусть дифференциальное уравнение приведено к виду $x = F(y, y')$, где
$F \in C^1$. Вводя параметр $p = y'$, получаем систему:
\begin{equation} \label{eq:l9_res_x}
    \begin{cases}
        x = F(y, p), \\
        dy = p \, dx.
    \end{cases} \tag{10}
\end{equation}

Продифференцируем первое уравнение системы \eqref{eq:l9_res_x}:
\[
    dx = \frac{\partial F}{\partial y} \, dy +
    \frac{\partial F}{\partial p} \, dp.
\]
Выразим из соотношения $dy = p \, dx$ дифференциал $dx = \dfrac{dy}{p}$ (в
предположении $p \neq 0$) и подставим его в левую часть:
\begin{equation} \label{eq:l9_res_x_diff}
    \frac{dy}{p} = \frac{\partial F}{\partial y}(y, p) \, dy +
    \frac{\partial F}{\partial p}(y, p) \, dp. \tag{11}
\end{equation}

\begin{enumerate}
    \item Интегрируя уравнение \eqref{eq:l9_res_x_diff} относительно $y$ и $p$,
    находим функцию $y = y(p, C)$ и записываем общее решение в параметрическом
    виде:
    \begin{equation} \label{eq:l9_res_x_ans}
        \begin{cases}
            x = F\bigl(y(p, C), p\bigr), \\
            y = y(p, C).
        \end{cases} \tag{12}
    \end{equation}
    \item \textbf{Важное правило проверки $p = 0$:} Поскольку при переходе к
    уравнению \eqref{eq:l9_res_x_diff} мы делили на $p$ ($dx = \frac{dy}{p}$),
    необходимо \textbf{всегда отдельно проверять случай $p = 0$}. Подставляя
    $p = 0$ в исходное уравнение, получаем соотношение $x = F(y, 0)$ (и
    проверяем, определяет ли условие $y = C, \; y' = 0$ решение исходного
    уравнения).
\end{enumerate}

\subsection{Примеры решения задач}

\begin{problem}[Ф.267]
    Решить дифференциальное уравнение:
    \begin{equation} \label{eq:l9_f267}
        x = (y')^3 + y'. \tag{13}
    \end{equation}
\end{problem}

\begin{proof}[Решение]~\\
    \begin{quote}
        Уравнение разрешено относительно независимой переменной $x$ и имеет вид
        $x = F(y')$ (частный случай уравнения $x = F(y, y')$). Введём параметр
        $p = y'$, тогда $dy = p \, dx$, а само уравнение принимает вид:
        \[
            x = p^3 + p.
        \]
        Найдём дифференциал $dx$:
        \[
            dx = \left(3 p^2 + 1\right) dp.
        \]
        Подставим выражение $dx = \dfrac{dy}{p}$ (при $p \neq 0$):
        \[
            \frac{dy}{p} = \left(3 p^2 + 1\right) dp \implies
            dy = \left(3 p^3 + p\right) dp.
        \]
        Интегрируя обе части по переменной $p$, находим $y(p)$:
        \[
            y = \int \left(3 p^3 + p\right) dp =
            \frac{3}{4} p^4 + \frac{1}{2} p^2 + C.
        \]
        Проверим отдельно случай $p = 0$: при $p = y' = 0$ имеем $y = C$, но
        подстановка $y' = 0$ в исходное уравнение \eqref{eq:l9_f267} даёт
        вертикальную прямую $x = 0$, которая не является графиком дифференцируемой
        функции $y(x)$ (а точка $(0, C)$ с горизонтальной касательной уже
        входит в параметрическое семейство при $p = 0$).

        \textbf{Ответ:}
        \[
            \begin{cases}
                x = p^3 + p, \\[1mm]
                y = \dfrac{3}{4} p^4 + \dfrac{1}{2} p^2 + C.
            \end{cases}
        \]
    \end{quote}
\end{proof}

\begin{problem}[Ф.272]
    Решить дифференциальное уравнение:
    \begin{equation} \label{eq:l9_f272}
        y = \ln\left(1 + (y')^2\right). \tag{14}
    \end{equation}
\end{problem}

\begin{proof}[Решение]~\\
    \begin{quote}
        Уравнение разрешено относительно $y$. Введём параметр $p = y'$, тогда
        $dy = p \, dx$, и уравнение записывается как:
        \[
            y = \ln\left(1 + p^2\right).
        \]
        Продифференцируем обе части:
        \[
            dy = \frac{2 p}{1 + p^2} \, dp.
        \]
        Заменим $dy$ на $p \, dx$ и перенесём все слагаемые в левую часть:
        \[
            p \, dx = \frac{2 p}{1 + p^2} \, dp \implies
            p \left(dx - \frac{2 \, dp}{1 + p^2}\right) = 0.
        \]
        Произведение обращается в нуль в двух случаях:
        \begin{enumerate}
            \item $dx = \dfrac{2 \, dp}{1 + p^2}$. Интегрируя, находим
            \[
                x = 2 \arctan p + C,
            \]
            что вместе с уравнением $y = \ln(1 + p^2)$ даёт общее решение в
            параметрической форме (при желании параметр $p$ можно исключить:
            $p = \tan\frac{x - C}{2} \implies y = -\ln\cos^2\frac{x - C}{2}$).
            \item $p = 0$. Подставляя $p = 0$ в выражение для $y$, получаем
            стационарное решение:
            \[
                y = \ln(1 + 0) = 0.
            \]
            Прямая $y = 0$ не получается из параметрического семейства ни при
            каком фиксированном $C$, но в каждой своей точке $(C, 0)$ (при
            значении параметра $p = 0$) касается соответствующей кривой
            семейства, то есть является \textdef{огибающей} (особым решением).
        \end{enumerate}

        \textbf{Ответ:}
        \[
            \begin{cases}
                x = 2 \arctan p + C, \\
                y = \ln\left(1 + p^2\right),
            \end{cases}
            \quad \text{а также особое решение } y = 0.
        \]
    \end{quote}
\end{proof}